The upshot is that we may now view $\FHG$ as an $\RGamma$-algebra, not merely an $\mathcal{R}$-algebra. But the fact that our ground ring now incorporates all conjugacy classes of $\Gamma$ means that rather than considering partially reduced cycle types, we will have to reduce all parts of a multipartition.
\begin{definition}
If an element of $\Gamma \wr S_n$ has cycle type $\bs \mu$, we say it has \emph{fully-reduced cycle type} $\bs \nu$, where for each $c \in \Gamma_*$, $\bs \nu(c)$ is obtained from $\bs \mu(c)$ by subtracting $1$ from each part, and ignoring any resulting parts of size zero. We let $\bm\hat{\bs \mu}$ be the multipartition obtained from $\bs \mu$ by adding 1 to each nonzero part of $\bs \mu(c)$ for each $c \neq 1$, and leaving $\bs \mu(1)$ unchanged.
\end{definition}

For an element of $\Gamma \wr S_n$, passing from the cycle type to the fully-reduced cycle type is generally not reversible, even if $n$ is known. While we may recover the number of $1$-cycles, their types may be arbitrary conjugacy classes. The notation has been set up so that $X_{\bm\hat{\bs \mu}}$ is the sum of all elements of $\Gamma \wr S_n$ with all 1-cycles labelled by the trivial conjugacy class and having fully-reduced cycle type $\bs \mu$.

\begin{theorem}
As $\RGamma$-module, $\FHG$ is free with basis $K_{\bm \hat{\bs \mu}}$ where $\mu \in \mathcal{P}(\Gamma_*)$.
\end{theorem}
\begin{proof}
By Proposition \ref{rgamma_free_over_r_prop}, $\RGamma$ is free over $\mathcal{R}$ with basis $B_{\mathbf{M}}$ indexed by $\mathbf{M}$ with $\mathbf{M}(1) = 0$. So it is enough to show that $K_{\bm\hat{\bs\mu}}B_{\mathbf{M}}$ is a $\mathcal{R}$-basis of $\FHG$. To accomplish this, we introduce a filtration of $\FHG$ by $\mathcal{R}$-modules by placing $K_{\bs \mu}$ in filtration degree $|\bs \mu| + l(\bs\mu(1))$. We now show that this respects multiplication, by passing to $Z(\mathbb{Z} \Gamma \wr S_n)$.

Let us say that $(\mathbf{g}, \sigma) \in \Gamma \wr S_n$ \emph{affects} $i \in \{1,\ldots, n\}$ if either $\sigma(i) \neq i$, or $\mathbf{g}_i \neq 1$. The number of elements of $\{1,\ldots,n\}$ that are affected by an element of partially-reduced cycle type $\bs \mu$ is precisely $|\bs \mu| + l(\bs\mu(1))$. Consider two elements $(\mathbf{g}, \sigma), (\mathbf{h}, \rho) \in \Gamma \wr S_n$, and suppose that their product $(\mathbf{g} \sigma(\mathbf{h}), \sigma \rho)$ affects $i$. Then either $\rho(\sigma(i)) \neq i$, or $\mathbf{g}_i\sigma(\mathbf{h})_i \neq 1$. If $\rho(\sigma(i)) \neq i$, then either $\sigma(i) \neq i$, or $\rho(i) \neq i$, which implies that one of $(\mathbf{g}, \sigma), (\mathbf{h}, \rho)$ affects $i$. If instead $\mathbf{g}_i\sigma(\mathbf{h})_i \neq 1$, either $\mathbf{g}_i \neq 1$ or $\sigma(\mathbf{h})_i \neq 1$. In the first of these cases, $(\mathbf{g}, \sigma)$ affects $i$. In the second case, either $\sigma(i) \neq i$, or $\mathbf{h}_i \neq 1$. Hence we conclude that if $(\mathbf{g} \sigma(\mathbf{h}), \sigma \rho)$ affects $i$, then one of $(\mathbf{g}, \sigma)$ and $(\mathbf{h}, \rho)$ affects $i$. This implies that the set of elements in $\{1,\ldots,n\}$ affected by the product of two group elements is a subset of the union of the elements affected by each factor. This proves that we have a filtration, and shows that only products of elements affecting disjoint subsets of $\{1,\ldots,n\}$ contribute to the leading term in the associated graded ring.

To compute $K_{\bm\hat{\bs\mu}}B_{\mathbf{M}}$, we again pass to $\mathbb{Z}\Gamma \wr S_n$. This gives $X_{\bm\hat{\bs\mu}}b_{\mathbf{M}}$. The leading order term arises from $X_{\bm\hat{\bs\mu}}$ and $b_{\mathbf{M}}$ affecting disjoint subsets of $\{1,\ldots, n\}$. In that case, we obtain $X_{\bs \nu}$, where $\bs \nu(c) = \bm \hat{\bs \mu}(c) \cup (1^{\mathbf{M}(c)})$ for $c \in \Gamma_*$. Any multipartition $\bs \nu$ arises from exactly one pair $(\bm\hat{\bs\mu}, \mathbf{M})$ in this way. In particular $\bs \mu$ is the fully-reduced cycle type corresponding to the partially-reduced cycle type $\bs \nu$, while $\mathbf{M}$ records how many parts of $\bs \nu$ of size $1$ have a given label in $\Gamma_*\backslash\{1\}$. We conclude that $K_{\bm\hat{\bs\mu}}B_{\mathbf{M}}$ equals $K_{\bs \nu}$ plus lower order terms. In particular the $K_{\bm\hat{\bs\mu}}B_{\mathbf{M}}$ form an $\mathcal{R}$-basis of $\FHG$ and the theorem~follows.
\end{proof}
